\subsection{Resolution of the two classes on $X_1(2,26)$}
\begin{proposition}\label{prop:decic226-complete}
Each of the two classes in Proposition~\ref{prop:decic226-residual}
has exactly one effective rational representative. Both representatives
are irreducible divisors of degree $10$, and their elliptic curves have
full torsion $C_2\times C_{26}$. Up to level automorphisms, these are
all noncuspidal points on $X_1(2,26)$ defined over decic fields.
In particular, $C_2\times C_{26}\in\Phi(10)$.
\end{proposition}
\begin{proof}
We give explicit witnesses. Let $K_i=\Q(a_i)$, where $a_i$ is a root of
\[
\begin{aligned}
 F_1(x)={}&2x^{10}+11x^9+17x^8-7x^7-42x^6-20x^5\\
          &\qquad+9x^4+9x^3+10x^2+6x+1,\\
 F_2(x)={}&x^{10}+3x^9-3x^8-23x^7-26x^6+3x^5\\
          &\qquad+13x^4-2x^3-9x^2-5x-1.
\end{aligned}
\]
Both polynomials are irreducible. Define
\[
 b_1=\frac{B_1(a_1)}{406832},\quad
 c_1=\frac{C_1(a_1)}{203416},\qquad
 b_2=\frac{B_2(a_2)}{784},\quad
 c_2=\frac{C_2(a_2)}{1960},
\]
where the coefficient table below specifies the four polynomials.
\begin{center}
\small
\begin{tabular}{r|rr|rr}
$j$&$[x^j]B_1$&$[x^j]C_1$&$[x^j]B_2$&$[x^j]C_2$\\\hline
0&-621631&129155&-27507&90619\\
1&-2586377&551533&-95105&307177\\
2&-2555940&1269576&-100181&325777\\
3&-4632176&2528052&101187&-321491\\
4&2207386&1757370&190643&-622235\\
5&13819447&-2078127&-236919&755143\\
6&4462921&-1932589&-349828&1128928\\
7&-4689287&-133261&-76640&249316\\
8&-3688261&347217&41312&-133040\\
9&-721870&92982&17265&-55877\\
\end{tabular}
\end{center}

Put
\[
 E_i:y^2=x^3+c_ix^2+b_i(1-b_i-c_i)x,\qquad
 P_i=(0,0),\quad Q_i=(b_i,b_i).
\]
Exact arithmetic gives $\operatorname{ord}(P_i)=2$,
$\operatorname{ord}(Q_i)=26$, and $P_i\ne13Q_i$.
Thus $E_i(K_i)$ contains $C_2\times C_{26}$.
The following good reductions bound the full torsion:
\[
\begin{array}{c|r|r|r|r|r}
 i&p&a_i\bmod\mathfrak p&b_i\bmod\mathfrak p&c_i\bmod\mathfrak p
       &\#E_i(\mathbb F_p)\\\hline
1&41&23&6&25&52\\
2&43&39&21&11&52.
\end{array}
\]
The defining polynomials are separable modulo the respective primes,
and the curve discriminants are nonzero there. The corresponding
local fields are unramified. Torsion injectivity at an odd unramified
prime therefore proves
\[
 E_i(K_i)_{\mathrm{tors}}\simeq C_2\times C_{26}.
\]

Let $D_i$ be the trace divisor of $(E_i,P_i,Q_i)$ on $X_1(2,26)$.
Reducing the defining polynomials at $3$ gives factor degrees $(5,5)$
for $K_1$ and $(10)$ for $K_2$. We evaluate $b_i,c_i$ at every root
in the corresponding finite fields. Their Frobenius orbits agree
exactly with the saved patterns defining $L_1,L_2$, respectively.
This checks the level points themselves, not only their $j$-invariants.
Consequently $[D_i-L_i]$ is rational torsion and reduces to zero at $3$.
Injectivity gives $D_i\sim L_i$ over $\Q$.

The $\mathbb F_3$-gonality is $12$, so each effective degree-ten
residue divisor has a one-dimensional Riemann--Roch space.
Semicontinuity now gives
\[
 1\leq h^0(X_1(2,26),L_i)\leq1.
\]
Hence $D_i$ is the unique effective representative of $L_i$.
It is irreducible of degree $10$: exact linear algebra verifies
$\Q(b_i,c_i)=K_i$. The native extraction also gives a single place
of degree $10$ in each case. Proposition~\ref{prop:decic226-residual} proves completeness.
The two $j$-invariants have different irreducible minimal polynomials,
both of degree $5$, so they also give distinct Galois orbits of elliptic
curves.
\end{proof}

These points are sporadic. Indeed, rank zero and gonality greater
than $10$ imply that $X_1(2,26)$ has only finitely many rational
effective divisors of degree $10$. The equations and all exact checks
are stored in \path{C226_resolution_20261030/}. In particular,
\path{complete_resolution_inputs.m} and the completed native and Sage
replays verify both witnesses without repeating the slow
characteristic-zero Riemann--Roch searches used to find them.

\begin{proposition}\label{prop:decic52}
No elliptic curve over a decic number field has a point of order $52$.
\end{proposition}
\begin{proof}
By Lemma~\ref{lem:decic-isogeny-transfer}, such a curve would be
$2$-isogenous to a curve containing $C_2\times C_{26}$.
Proposition~\ref{prop:decic226-complete} identifies this curve, after
field conjugacy and a change of level basis, with one of the two
curves above. Its level point has degree $10$, so its field of
definition is the entire decic field.

Write the roots of its $2$-division polynomial as
\[
 r_0=0,\qquad r_1=x(13Q_i),\qquad r_2=-c_i-r_1.
\]
For a kernel generated by $(r_s,0)$, translate $r_s$ to zero and put
$u=r_t-r_s$, $v=r_w-r_s$, where $\{s,t,w\}=\{0,1,2\}$.
The quotient has equation
\[
 E':y^2=x^3+2(u+v)x^2+(u-v)^2x.
\]
The dual kernel is generated by $(0,0)$.
This point is divisible by $2$ over $K_i$ if and only if $u$ or $v$
is a square in $K_i$. Indeed, the duplication formula forces the
$x$-coordinate of a half-point to be $u-v$ or $v-u$; its equation
then requires a square root of $u$ or $v$.

None of the six differences $r_t-r_s$, $s\ne t$, is a square.
Here are exact residue certificates. Each displayed value is a
nonzero nonsquare in the indicated prime field.
\[
\begin{array}{c|r|rr}
 (s,t)&K_1:\ p=41,\ a_1=23
       &\multicolumn{2}{c}{K_2:\ p=43}\\
 &r_t-r_s&a_2&r_t-r_s\\\hline
 (0,1)&38&16&33\\
 (0,2)&19&39&18\\
 (1,0)&3&39&29\\
 (1,2)&22&16&19\\
 (2,0)&22&16&34\\
 (2,1)&19&39&39.
\end{array}
\]
The field equations vanish at these residues and are separable
modulo the primes. A square in $K_i$ which is a unit at such a place
would have square residue, proving the assertion.

Finally, in the isogeny-transfer construction from a point $R$ of
order $52$, the point $13R$ has order $4$ and doubles to the original
isogeny kernel $\langle26R\rangle$. Under the dual quotient this is
exactly the kernel point $(0,0)$ considered above. Its required
half-point does not exist for either curve and any of the three
kernels. This is a contradiction.
\end{proof}
